\subsection{Completeness}\label{sec:completeness}

Let us now turn our attention to proving the completeness of tableaux as a calculus for satisfiability.
The main idea behind the proof is to look at an open and saturated branch of tableaux and to build a model from it (cf., e.g., \cite{Fit83,Fit96,Priest08}).
%To simplify the presentation of proofs, in this section we assume $\noms$ is the set $\mathbb{N}$  of natural numbers.
We begin with some preliminary definitions and results.

% \begin{defi}
%     We use $\noms(\Theta)$ for the set of nominals appearing in a branch $\Theta$ of a tableau.
% \end{defi}

\begin{defi}
    \label{def:prefixrelation}
    We use $\noms(\Theta)$ for the set of nominals appearing in a branch $\Theta$ of a tableau. 
    If $\Theta$ is saturated, for every $\set{i,j} \subseteq \nominal(\Theta)$ define
    $i \equiv_\Theta j \mbox{ iff } i{:}j \in \Theta$, and let $[i]_{\equiv_{\Theta}} = \{ j \mid i \equiv_\Theta j\}$.
\end{defi}

Intuitively, the relation $\equiv_\Theta$ groups together all the nominals in the branch $\Theta$ representing the same node in the model we want to build.
The following lemma states $\equiv_\Theta$ is an equivalence relation as desired.

\begin{lem}
    \label{lemma:equiv-prefixrelation}
    $\equiv_\Theta$ is an equivalence relation.
\end{lem}
\begin{proof}
Immediate from (ref), (sym) and (trans).
%Transitividad sale de usar Ref, Sym, y Nom. Si tenemos ab y bc, usamos que de ba bb cb implica ca
\end{proof}

Let us proceed to introduce key elements in our proof completeness: \emph{nominal urfathers}.

\begin{defi}\label{def:urfather} 
    Let $\Theta$ be a saturated tableau branch of a tableau.
    Define $\urf[i]=\min([i]_{\equiv_{\Theta}})$; and we say that
$\urf[i]$ is the \emph{urfather} of $i$.
    %And if $\urf[i] = i$, then, $i$ is a \emph{nominal urfather}. 
    The set $\urfather{}{\Theta} = \setof{\urf[i]}{i \in \noms(\Theta)}$ is the set of urfathers in the branch.
\end{defi}

The following is an important property of nominal urfathers.
Intuitively, as pointed out in \cite{Brauner2011}, this property tells us that non-trivial reasoning involving non-singleton equivalence classes is carried out only in connection with nominals appearing in the root node expression.

\begin{lem}\label{lemma:representative}
    If $\Theta$ is a saturated branch of a tableau and $\urf \neq i$, $\urf$ is a root nominal.
\end{lem}
\begin{proof}
    Let $\urf \neq i$ and $\urf = j$.
    The proof proceeds by cases.
    
    \begin{description}
        \item[($i$ is a root nominal)]
        Recall that a root nominal is a nominal appearing in the root of our tableau. Any new nominal added by the application of a rule from~\Cref{subsec:rules} is greater than any root nominal. Since $j < i$ and $i$ is a root nominal, $j$ must be also a root nominal.
        % Immediately, $j$ must be a root nominal (otherwise, we would get a contradiction, since we would have $j > i$, as new nominals are greater than root nominals, and $j < i$).

        \item[($i$ is not a root nominal)]
        If $i{:}j$ is the root of the tableau, then, we are done.
        If $i{:}j$ is not the root of the tableau, then, it was introduced by a tableau expansion rule.
        It is clear that $j$ is a root nominal if it was introduced by ($\lnot\lnot$), ($\land$), (nom), and (test).
        Notice that in any of these cases, $j$ is a component of the node expression labelling the root of the tableau.
        Otherwise, $i{:}j$ is introduced into the tableau by (sym) or (trans).
        In these two last cases, it must be that there is a root nominal $k$ such that $\set{i{:}k, k{:}k} \subset \Theta$.
        That $j$ is a root nominal follows from this fact and the minimality of being an urfather.
    \end{description}
\end{proof}

% \begin{defi} \label{def:urfather:2}
%     {\color{red} Let $\Theta$ be a saturated branch of an tableau, and let $i$ be any nominal occurring on  $\Theta$. The \emph{urfather} of $n$ on  $\Theta$, denoted $\uf(n)$, is defined by
%     \[
%         \uf(n)=\min([n]_{\Theta}) 
%     \]  
%     } 
% \end{defi}
From their definition, it follows that nominal urfathers can be taken as canonical representatives of the equivalence classes in $\noms(\Theta)/{\equiv\Theta}$.
This characteristic makes them key elements in our proof of completeness.
Moreover, as made clear in \Cref{lemma:urfclos}, nominal urfathers collect the information about all members of the equivalence class.

% , crucial in the construction of a model from an open and saturated branch of a tableau,

% The proof of~\Cref{lemma:urfclos} makes clear that this is the goal of the rule (copy) in~\Cref{rules:nominals}.

\begin{lem}\label{lemma:urfclos}
    If $\Theta$ is a saturated branch of a tableau and $i{:}\varphi \in \Theta$, then, $\urf{:}\varphi \in \Theta$.
\end{lem}
\begin{proof}
    Suppose that  $i{:}\varphi\in\Theta$ and $\urf=j$.
    If $i = j$, then, we are done.
    If $i \neq j$, from \Cref{lemma:representative}, we know $j$ is a root nominal.
    Moreover, from \Cref{def:prefixrelation},  $i{:}j \in \Theta$, and, from \Cref{def:urfather}, we know $j < i$.
    Since $\Theta$ is saturated it is closed under (copy), $j{:}\varphi\in\Theta$.
    Thus, $\urf:\varphi \in \Theta$.
\end{proof}

The results in \Cref{lemma:representative,lemma:urfclos} pave the way to define the intended model from an open branch.
Intuitively, they tell us that if our goal is to build a model that satisfies a node expression $\varphi$ from an open and saturated branch of a tableau, then, we can restrict our attention to nominal urfathers.
In comparison with what is done in~\cite{ArecesFS17}, the use of urfathers simplifies the definition of the model, and aid in the proof of \Cref{prop:truthlemma}.

\begin{defi}[Extracted Model]\label{def:extmodel}
    Let $\Theta$ be an open and saturated branch of a tableau.
    Define a structure
        $ \model^\Theta =
            \tup{
                N^{\Theta},
                \set{R^{\Theta}_\pathh}_{\pathh \in \rels},
                \set{\approx^{\Theta}_{\cmp}}_{\cmp\in\cmps},
                \valuation^{\Theta},
                \naming^{\Theta}
            }$
    where
    \[
        \begin{array}{l@{\ }c@{\ }l}
            N^{\Theta}              & = & \urfather{}{\Theta} \\
            R^{\Theta}_\pathh       & = & \{(\urf,\urfather{\Theta}{j}) \mid  i{:}\tup{\child}j\in\Theta \} \\
            \approx^{\Theta}_{\cmp} & = & \{(\urf, \urfather{\Theta}{j}) \mid \tup{@_i =_\cmp @_{j}} \in\Theta \}\\            
            \valuation^{\Theta}(p) & = & \{\urf \mid  i{:}p \in \Theta \} \\
            \naming^{\Theta}(i) & = & 
                    \begin{cases}
                        \min(\urfather{}{\Theta}) & \text{if }  i \notin\nominal(\Theta)\\
                        \urfather{\Theta}{i} & \text{if } i\in\nominal(\Theta).
                    \end{cases} 	
        \end{array}
    \]
    % \begin{align*}
    %     N^{\Theta} & = \urfather{}{\Theta} \\
    %     R^{\Theta}_\pathh & = \setof{(\urf,\urfather{\Theta}{j})}{i{:}\tup{\child}j\in\Theta} &
    %     \valuation^{\Theta}(p) & = \setof{\urf}{i{:}p \in \Theta} \\
    %     \approx^{\Theta}_{\cmp} & = \setof{(\urf, \urfather{\Theta}{j})}{\tup{@_i =_\cmp @_{j}} \in\Theta} &
    %     \naming^{\Theta}(i) & = 
    %         \urfather{\Theta}{i} \text{ if } i\in\nominal(\Theta) \text{ else } \min(\urfather{}{\Theta})
    % \end{align*}
\end{defi}

\begin{prop}\label{prop:extmodel}
    The structure $\model^{\Theta}$ in \Cref{def:extmodel} is a model of $\Hxpath$.
\end{prop}
\begin{proof}
    The only challenge is proving $\approx^{\Theta}_{\cmp}$ to be an equivalence relation. 
    Reflexivity follows from (ref) and (dRef), symmetry follows from (com$_1$), and transitivity follows from (dTrans).
\end{proof}

We refer to $\model^{\Theta}$ as the \emph{extracted model} for $\Theta$.
This is the main ingredient in the proof of completeness of tableaux as a calculus for satisfiability.
The proof of this result is carried out by induction on the \emph{size} of a node expression. 
We introduce this definition of size below and comment on it immediately after. 

\begin{defi}\label{def:size}
    The function $\size$ is defined by mutual recursion on node and path expressions.
    \[\begin{array}{l@{\,}l@{\qquad}@{\qquad}l@{\,}l}
        \size(i) & = 1 & %\text{~~~for $i\in\noms$} &
        \size(\child ) & = 1 \\ % \text{~~~for $\child\in\mods$} \\
        \size(p) & = 1 & % \text{ ~~for $p\in\props$} &
        \size(@_i) & = 1 \\ % \text{~~~for $i\in\noms$} \\
        \size(\varphi \land \psi) & = 1 + \size(\varphi) + \size(\psi) &
        \size(\varphi?) & = 1 + \size(\varphi) \\
        \size(\lnot \varphi) & = 1 + \size(\varphi) &
        \size( \alpha \beta ) & = \size(\alpha) + \size(\beta) \\
        \size(\tup{\child} \varphi) & = 1 + \size(\varphi) &
        \size( \alpha \cup \beta ) & = 1 + \size(\alpha) + \size(\beta) \\
        \size(i{:}\varphi) & = 3 + \size(\varphi) \\
        \size(\tup{\alpha \cmpr \beta}) & = 5 + \size(\alpha) + \size(\beta).
    \end{array}\]   
\end{defi}

The definition of $\size$ is necessary because we cannot rely on structural induction over node expressions. This is due to the fact that rule applications may generate expressions not structurally contained in the original input. For example, applying ($\lnot$int) to $i{:}\lnot\tup{\alpha \cmpr \beta}$ introduces $\lnot\tup{@_i\alpha \cmpr @_i\beta}$ in the branch, which is not a part of the original expression.
To handle this, we define a size measure on node and path expressions and carry out induction on that measure instead. The constants $3$ for $i{:}\varphi$ and $5$ for $\tup{\alpha \cmpr \beta}$ are chosen to ensure that rule applications strictly reduce size. This guarantees that the induction is well-founded. For instance, in the case of ($\lnot$int), we have $\size(\lnot\tup{@_i\alpha \cmpr @_i\beta}) < \size(i{:}\lnot\tup{\alpha \cmpr \beta})$.

\begin{lem}\label{prop:truthlemma}
    Let $\Theta$ be an open and saturated branch, and $\model^{\Theta}$ be the model extracted from it. 
    In addition, let $\model^{\Theta} \vDash \varphi$ indicate $\model^{\Theta}, n \vDash \varphi$ for all $N^{\Theta}$.
    Then, $\varphi \in \Theta$ implies $\model^{\Theta} \vDash \varphi$.
\end{lem}
\begin{proof}
    Observe that every $\varphi \in \Theta$ has one of the following forms: 
    \[
        (1)~i{:}\psi \qquad\qquad
        (2)~\tup{@_i \alpha \cmpr @_j\beta } \qquad\qquad
        (3)~\lnot\tup{@_i \alpha \cmpr @_j\beta }
    \]
    The proof is by induction on the size of $\varphi \in \Theta$.
    We focus on some of the cases in the proof.
    \begin{description}
        \item[\textnormal{For the base cases we consider}] \hfill
            \begin{description}
                \item[($\varphi = i{:}j$)] 
                    % Let $i{:}j \in \Theta$.
                    We need to prove $\model^\Theta \vDash i{:}j$, or equivalently, $\urf[i] = \urf[j]$.
                    The result is immediate from the definition of $\naming^{\Theta}$.
                \item[($\varphi = i{:}\lnot j$)]
                    % Let $i{:}\lnot j \in \Theta$.
                    We need to prove $\model^\Theta \vDash i{:}\lnot j$, or equivalently $\urf \neq \urfather{\Theta}{j}$.
                    The proof is by contradiction.
                    Suppose $\urf = \urf[j]$.
                    By \Cref{def:urfather,def:prefixrelation}, $i{:}j \in \Theta$.
                    This implies that $\Theta$ contains a clash, contradicting that $\Theta$ is open. Thus, $\model^\Theta \vDash i{:}\lnot j$. 
                \item[($\varphi = \tup{@_i =_{\cmp} @_j}$)] 
                    % Suppose $\tup{@_i =_{\cmp} @_j} \in \Theta$.
                    We neet to prove $\model^\Theta \vDash \tup{@_i =_{\cmp} @_j}$, or equivalently that the pair ${(\urfather{\Theta}{i}, \urfather{\Theta}{j})} \in {\approx^{\Theta}_{\cmp}}$.
                    The result is immediate from the definition of ${\approx^{\Theta}_{\cmp}}$.
                \item[($\varphi = \tup{ @_i \neq_{\cmp} @_j}$)] 
                    % Suppose $\tup{@_i \neq_{\cmp} @_j} \in \Theta$.
                    We need to prove $\model^\Theta \vDash \tup{@_i \neq_{\cmp} @_j}$, or equivalently that the pair ${(\urf, \urfather{\Theta}{j})} \notin {\approx^{\Theta}_{\cmp}}$.
                    The proof is by contradiction.
                    Suppose that the pair ${(\urf, \urfather{\Theta}{j})} \in {\approx^{\Theta}_{\cmp}}$.
                    By definition of ${\approx^{\Theta}_{\cmp}}$, $\tup{@_i =_{\cmp} @_j} \in \Theta$.
                    But this means the branch contains a clash, contradicting that $\Theta$ is open. Thus, $\model^\Theta \vDash \tup{@_i \neq_{\cmp} @_j}$. 
            \end{description}
        \item[\textnormal{The inductive hypothesis (IH) is}] for all $\psi \in \Theta$ such that $\size(\psi) < \size(\varphi)$,  $\model^\Theta \vDash \psi$.
        \item[\textnormal{For the inductive cases we consider}] \hfill
        \begin{description}
            % \item[($\varphi = i{:}(\psi \land \chi)$)] 
            %     Suppose $i{:}(\psi \land \chi) \in \Theta$.
            %     We need to prove $\model^\Theta \vDash i{:}(\psi \land \chi)$.
            %     Since $\Theta$ is saturated, we know $\set{i{:}\psi, i{:}\chi} \subseteq \Theta$.
            %     The result follows immediately from the IH.
            % \item[($\varphi = i{:}\lnot \psi$)]
            %     It follows that $\psi$ is either one of:
            %     \begin{description}  
            \item[($\varphi = i{:}\lnot\lnot \psi$)]
                We need to prove $\model^{\Theta} \vDash i{:}\lnot\lnot\psi$, or equivalently
                    $\model^{\Theta} \vDash i{:}\psi$.
                From $(\lnot\lnot)$, ${i{:}\psi \in \Theta}$.
                By the IH, $\model^{\Theta} \vDash i{:}\psi$.
            \item[($\varphi = i{:}\lnot (\psi \land \chi)$)]
                We need to prove $\model^{\Theta} \vDash i{:}\lnot(\psi \land \chi)$, or equivalently
                    $\model^{\Theta} \vDash i{:}\lnot\psi$ or
                    $\model^{\Theta} \vDash i{:}\lnot\chi$.
                From ($\lnot\land$), $i{:}\lnot\psi \in \Theta$ or $i{:}\lnot\chi \in \Theta$.
                W.l.o.g., let $i{:}\lnot\psi \in \Theta$. By the IH, $\model^{\Theta} \vDash i{:}\lnot\psi$.
            \item[($\varphi = i{:}\lnot j{:}\psi$)]
                We need to prove $\model^{\Theta} \vDash {i{:}\lnot(j{:}\psi)}$, or equivalently
                    $\model^{\Theta} \vDash {j{:}\lnot\psi}$.
                From ($\lnot$nom), $j{:}\lnot\psi \in \Theta$.
                Then, by the IH, $\model^{\Theta} \vDash j{:}\lnot \psi$.
            \item[($\varphi = i{:}\lnot \tup{\pathh} \psi$)]
                We must prove $\model^{\Theta} \vDash {i{:}\lnot \tup{\pathh}\psi}$, or equivalently,
                for all $(\urf,n) \in R^{\Theta}_{\pathh}$, ${\model^{\Theta}, n \vDash \lnot\psi}$.
                Let $(\urf,n) \in R^{\Theta}_{\pathh}$.
                The definition of $R^{\Theta}_{\pathh}$ tells us that exists $i{:}\tup{\pathh}j \in \Theta$ s.t.\ $\urfather{\Theta}{j} = n$.
                Since $\{i{:}\lnot\tup{\pathh}\psi, i{:}\tup{\pathh}j\} \subseteq \Theta$, using ($\lnot \Diamond$), we get $j{:}\lnot\psi \in\Theta$.
                By IH, $\model^\Theta \vDash j{:}\lnot \psi$, i.e., $\model^\Theta, n \vDash \lnot \psi$.
            \item[($\varphi = i{:}\lnot\tup{\alpha \cmpr \beta}$)]
                We neet to prove $\model^{\Theta} \vDash i{:} \lnot \tup{\alpha \cmpr \beta}$.
                Equivalently, we prove $\model^{\Theta} \vDash {\lnot \tup{@_i\alpha \cmpr @_i\beta}}$.
                From (int$_2$), $\lnot\tup{ @_i \alpha \cmpr @_i \beta} \in \Theta$.
                Then, we have $\size(\lnot\tup{ @_i \alpha \cmpr @_i \beta}) < \size(i{:} \lnot \tup{\alpha \cmpr \beta})$.
                This enables us to apply the IH and get $\model^\Theta \vDash \lnot\tup{@_i \alpha \cmpr @_i \beta}$.
                % \end{description}
            % \item[($\varphi = i{:}j{:}\psi$)]
            %     Suppose $i{:}j{:}\psi \in \Theta$.
            %     We need to prove $\model^{\Theta} \vDash i{:}j{:}\psi$.
            %     From the definition of~$\vDash$, this is equivalent to proving $\model^{\Theta} \vDash j{:}\psi$.
            %     Since $\Theta$ is saturated, using (nom), we know $j{:}\psi \in \Theta$.
            %     Therefore, by IH, $\model^\Theta \vDash j{:}\psi$.
            \item[($\varphi = i{:}\tup{\child}\psi$)] 
                % Suppose $i{:}\lnot \tup{\child} \psi  \in \Theta$.
                We need to prove $\model^{\Theta} \vDash {i{:}\tup{\child}\psi}$.
                From ($\Diamond$), $\set{i{:}\tup{\child}j, j{:}\psi} \subseteq \Theta$. 
                Using the IH, we can prove $\model^\Theta \vDash i{:}\tup{\pathh}j \land j{:}\psi$.
                \Cref{lemma:child} gives us the desired result.
            %%%% TYPE 2 %%%%
            % \item[($\varphi = i{:}\tup{\alpha \cmpr \beta}$)] 
            %     Similar to $\varphi = i{:}\lnot\tup{\alpha \cmpr \beta}$ but using (int$_1$) instead of (int$_2$).
            % \item[($\varphi =  \tup{@_i \child \alpha \cmpr \beta}$)]
            %     Similar to $\varphi = i{:}\tup{\child}\psi$ but using (child) instead of ($\Diamond$), and~\Cref{lemma:child}, item (2). 
            % \item[($\varphi =  \tup{@_i@_j\alpha \cmpr \beta}$)]
            %     Similar to $\varphi = i{:}j{:}\psi$ but using ($@$) instead of (nom).
            % \item[($\varphi =  \tup{@_i\varphi?\alpha \cmpr \beta}$)]
            %     Similar to $\varphi = i{:}(\psi \land \chi)$ but using (test) instead of ($\land$).
            \item[($\varphi = \tup{@_i (\alpha \cup \beta) \gamma \cmpr \eta}$)]
                % Suppose $\tup{@_i (\alpha \cup \beta) \gamma \cmpr \eta} \in \Theta$.
                We need to prove that $\model^\Theta \vDash \tup{@_i (\alpha \cup \beta) \gamma \cmpr \eta}$, equivalently $\model^\Theta \vDash \tup{@_i\alpha\gamma \cmpr \eta}$ or $\model^\Theta \vDash \tup{@_i \beta\gamma \cmpr \eta}$.
                From ($\cup$), $\tup{@_i\alpha\gamma \cmpr \eta} \in \Theta$ or $\tup{@_i\beta\gamma \cmpr \eta} \in \Theta$.
                W.l.o.g.\ let $\tup{@_i\alpha\gamma \cmpr \eta} \in \Theta$, by the IH, $\model^{\Theta} \vDash \tup{@_i\alpha\gamma \cmpr \eta}$.
                % Therefore, $\model^{\Theta} \vDash \tup{@_i(\alpha \cup \beta)\gamma \cmpr \eta}$.
            % \item[($\varphi = \tup{@_i \cmpr \beta}$)]
            %     We can use (com$_1$) and think of the case analysis for $\beta$ in terms of the proofs for $\varphi = \tup{\beta \cmpr @_i}$.
            %%%% TYPE 3 %%%%
            % \item[($\varphi =  \lnot\tup{@_i\child \alpha \cmpr \beta}$)]
            %     Similar to $\varphi = i{:}\lnot\tup{\child}\psi$ but using ($\lnot$child) instead of (${\lnot}{\Diamond}$).
            % \item[($\varphi =  \lnot\tup{@_i@_j\alpha \cmpr \beta}$)]
            %     Similar to $\varphi = i{:}(\lnot j{:}\psi)$ but using (${\lnot}@$) instead of ($\lnot$nom).
            % \item[($\varphi =  \lnot\tup{@_i\varphi?\alpha \cmpr \beta}$)]
            %     Similar to $\varphi = i{:}\lnot(\psi \land \chi)$ but using ($\lnot$test) instead of (${\lnot}{\land}$).
            \item[($\varphi = \lnot\tup{@_i(\alpha \cup \beta) \gamma \cmpr \eta}$)]
                % Suppose $\lnot\tup{@_i (\alpha \cup \beta) \gamma \cmpr \eta} \in \Theta$.
                We need to prove  $\model^\Theta \vDash \lnot\tup{@_i (\alpha \cup \beta) \gamma \cmpr \eta}$, equivalently $\model^\Theta \vDash \lnot\tup{@_i\alpha\gamma \cmpr \eta}$ and $\model^\Theta \vDash \lnot\tup{@_i \beta\gamma \cmpr \eta}$.
                From (${\lnot}{\cup}$), $\lnot\tup{@_i\alpha\gamma \cmpr \eta} \in \Theta$ and $\lnot\tup{@_i\beta\gamma \cmpr \eta} \in \Theta$.
                By the IH, $\model^{\Theta} \vDash \lnot\tup{@_i\alpha\gamma \cmpr \eta}$ and $\model^{\Theta} \vDash \tup{@_i\beta\gamma \cmpr \eta}$.
            % \item[($\varphi = \lnot\tup{@_i \cmpr \beta}$)]
            %     We can use (com$_2$) and think of the case analysis for $\beta$ in terms of the proofs for $\varphi = \lnot\tup{\beta \cmpr @_i}$.
        \end{description}

    \end{description}
\end{proof}

The lemma above is the fundamental piece to state the intended completeness result.

\begin{thm}[Completeness]
    Let $\varphi$ be a node expression.
    If there is a tableau for $\varphi$ with an open and saturated branch, then, $\varphi$ is satisfiable.
\end{thm}
\begin{proof}
    Suppose that there is a tableau for $\varphi$ with an open and saturated branch $\Theta$.
    Let $\model^{\Theta}$ be the extracted model for $\Theta$.
    From \Cref{prop:truthlemma}, we know $\model^{\Theta} \vDash i{:}\varphi$ where $i{:}\varphi$ is the root of the tableau.
    Immediately, $\model^{\Theta}, \naming(i) \vDash \varphi$.
    Thus, $\varphi$ is satisfiable.
\end{proof}

As it is usual, our tableaux-based calculus not only provides a way to check the satisfiability of a node expression, but also acts as a \emph{model building procedure}, since for every satisfiable node expression it is possible to obtain a model (\Cref{def:extmodel}). 
We finish this section by showing how the extracted model is built for a given node expression.

\begin{exa}
    \label{ex:extracted model}
    Consider the tableau for $\varphi = \tup{\child}\tup{@_2\astep[b]2? =_\cmp \astep[b](q\land 3)?}$ in~\Cref{ex:usetableaux}.
    This tableau has a sole open and saturated branch $\Theta$.
    Thus, we can conclude $\varphi$ is satisfiable.
    We depict the extracted model from $\Theta$ in \Cref{ex:extracted:model}.
    In this figure, nodes are labelled by the nominals and propositional symbols satisfied by them.
    We omit the reflexive edges for $\approx_\cmp$.
    \begin{figure}[t]
        \centering
        \begin{tikzpicture}%
            [
                every node/.style={scale=.7},
                modal
            ]
            \node[world] (4) [label=225:{\footnotesize $4$}] {$4$};
            \node[world] (5) [label=225:{\footnotesize $5$},below=of 4] {$5$};
            \node[world] (7) [label=315:{\footnotesize $3$},below=of 5] {$3$, $7$, $q$};
            \node[world] (6) [label=225:{\footnotesize $2$},left=2cm of 7] {$2$, $6$};

            \path[->]
                (4)
                edge
                node[midway, right]{\footnotesize$\mathsf{a}$}
                (5);
            \path[->]
                (5)
                edge
                node[midway, right]{\footnotesize$\mathsf{b}$}
                (7);
            \path[-]
                (7)
                edge[dashed, bend right]
                node[midway, above]{\footnotesize$\approx_\cmp$}
                (6);
            \path[->]
                (6)
                edge[reflexive]
                node[midway, above]{\footnotesize$\mathsf{b}$}
                (6);
        \end{tikzpicture}
        \caption{Extracted model for $\tup{\child}\tup{@_2\astep[b]2? =_\cmp \astep[b](q\land 3)?}$}
        \label{ex:extracted:model}
    \end{figure}
    Notice that in the extracted model nodes correspond to urfathers; in this respect, we have $\urf[7]=3$, $\urf[6]=2$, and for any other nominal $k$, $\urf[k]=k$.
\end{exa}

Up to now we focused on node expressions.  Handling satisfiability of path expressions is simple, given the following proposition. 

\begin{prop}
	Let $\alpha$ be a path expression of $\Hxpath$, then $\alpha$ is satisfiable if and only if $i{:}\tup{\alpha}\top$ is satisfiable, for $i$ a nominal not in $\alpha$. 
\end{prop}	

\begin{proof}
	The implication from right to left is obvious.  For the other direction, any model satisfying $\alpha$ can be turn into a model of $i{:}\tup{\alpha}\top$ by setting $\naming(i)$ to the starting point of the path satisfying $\alpha$. 
\end{proof}	
	